% ===========================================================================
\section{Limitations}
\label{sec:limitations}
% ===========================================================================

\textbf{The sewer methane correction factor is unresolved, and it is the sector.} A factor of
\num{1.59} on two-thirds of the emissions, argued in Section~\ref{sec:discussion}. Both
readings are published; neither is measured.

\textbf{Industrial solid waste and industrial wastewater are not estimated at all.} No national
quantity for either has been compiled here. The categories reported here are municipal and
domestic only, and both gaps carry the notation key NE with that justification, in the
complete sector-4 matrix of \ref{app:completeness}---which also declares biological
treatment (4B) and incineration (4C1) as NE rather than leaving them silently absent.
Industrial wastewater is the more consequential omission: given that domestic wastewater
alone is \SI{5.5}{\percent} of the national inventory, an industrial stream from a textile
and tanning sector of Pakistan's size is unlikely to be negligible.

\textbf{Generation rates are single national values.} The urban and rural per-capita rates do
not vary by province or city size, which is a real simplification for a country whose largest
city and smallest districts differ substantially in consumption and in what they discard. The
rates are regional literature values rather than Pakistani measurements.

\textbf{The collection rate is the softest number in the partition.} The mass balance
closes by construction, but \SI{50.3}{\percent} is an assumption, and it governs how much
waste decays slowly in a landfill versus burns immediately in the street. It is a single
national figure applied through an urban/rural split, when collection in practice varies
between municipalities far more than between those two categories. The one Pakistani survey
that asks the question directly, HIES 2024--25, reports materially less collection than the
assumption (Section~\ref{sec:validation}); its household-route shares are priced as a
structural scenario rather than adopted, because route shares are not mass shares.

\textbf{The burned fraction is a literature value.} That \SI{60}{\percent} of uncollected
waste burns is the least evidenced parameter in the sector and has no Pakistani measurement
behind it at all. It scales open burning linearly, and with it the entire air-pollutant layer
in Section~\ref{sec:air}. Its complement is the fate ledger's unresolved row:
\SI{6.27}{\mega\tonne} of uncollected, unburned waste---dumped to land and
waterways---to which no emissions pathway is applied at all. That choice is now bounded
rather than silent: the scenario table prices the unresolved mass as unmanaged shallow
disposal (sector total \SI{39.87}{\mega\tonne}) and unmanaged deep disposal
(\SI{42.28}{\mega\tonne}) against the reference reading (\SI{37.45}{\mega\tonne}),
on the declared assumption that the dumped stream shares the collected stream's deposit
history and composition. Waterway decay lies between inert and shallow. No probability
attaches to any point in the bound---a disposal-site census, not a better default, is
what would close it.

\textbf{Waste composition is regional, not Pakistani.} The fossil carbon fraction governs how
much of the burning \COtwo{} is reportable, and the plastics and textiles content governs the
chlorine and much of the particulate. A Pakistani waste-characterisation survey would replace
both.

\textbf{No methane recovery is credited and no landfill oxidation applied.} Both are set to
zero because neither is documented at national scale, and both choices bias the disposal
category high if any large Pakistani site does in fact manage its gas or maintain a cover. The
direction of that bias is stated rather than hedged.

\textbf{The spin-up back-projection is an assumption, not data.} Extending the deposit series
to 1960 at \SI{40}{\percent} of the panel's first year is a judgement about Pakistan's waste
history. Its effect is reported---\SI{12}{\percent} of current landfill methane---so a reader
can remove it, but the shape of the back-projection is not evidenced.

\textbf{Uncertainties are large, and honestly so---and honestly asymmetric.}
Five of the six streams carry between
\SI{74}{\percent} and \SI{124}{\percent}; wastewater N$_2$O carries the bounded 150-fold
factor range, and the sector reads $\pm\SI{51.1}{\percent}$ by quadrature with a
parameter-rerun 95\% interval of \SIrange{25.7}{61.4}{\mega\tonne}
(Supplementary Note~S1).
These follow IPCC's own Tier~1 guidance for waste parameters and reflect that the underlying
activity is modelled from population rather than measured at facilities. They are not an
artefact of conservative parameter choice; they are what Tier~1 can say about a country with
no facility-level waste data.

\textbf{Dioxins and furans are not estimated.} Open burning of chlorinated plastics produces
them, and the hydrogen chloride reported in Section~\ref{sec:air} is a marker for the
conditions that do. Estimating them needs a dedicated study rather than a factor applied to
the same tonnage.
